\section{Proof of the Trilemma}
\label{app:proof}

Two formalisation choices come first.

\paragraph{Where the words enter.}
One might let each word enter only through its initialisation
$\Q_{z_i}^{(0)} \propto \exp(\Sfac_{w_i,\cdot})$. In the encoder model it also enters elsewhere:
$w_i$ appears in the unary term of $E$, so $\Sfac_{w_i,\cdot}$ reappears in every update of
$\Zv{i}$. We therefore assume only \emph{word-locality} (A3 below); the initialisation-only
version is the case $u_i \equiv 0$. The assumption matters. With word-nonlocal terms the theorem
is false, as
%
\begin{equation}
\label{eq:nonlocal-witness}
  E(\Q; w) = -\sum_i \sum_a \Q_i(a)\Big(\Sfac_{w_i,a}
     + \sum_{k<i} \Sfac'^{(i,k)}_{w_k,a}\Big)
\end{equation}
%
shows. This is a single smooth energy that is causal and contextual, because it conditions on
the prefix directly and bypasses inference. The theorem is about context mediated by inference
over the marginals, and (A3) states that scope formally.

\paragraph{What causality means.}
(C1) \emph{Trajectory causality}: for all $i<j$, all schedules and all $t$, the iterate
$\Q_{z_i}^{(t)}$, as a function of the words, does not depend on $w_j$. (C2) \emph{Network
causality}: for every schedule, every $t$ and every node $x$, head nodes included,
$\Q_x^{(t)}$, as a function of an arbitrary interior initialisation, does not depend on the
initialisation blocks of nodes at strictly later positions. The theorem is false under (C1). A
finite vocabulary makes (C1) a finite family of conditions on the states the trajectories
actually visit, whereas the graph condition the proof needs is an identical-vanishing condition
on the whole interior of the state space.

\begin{proposition}[Sharpness]
\label{app:prop-sharp}
Some scheme satisfies SE, C1 and X simultaneously.
\end{proposition}

\begin{proof}
Take $\seqlen=2$, binary labels ($\nlab=2$) and no head variables (permitted: the factor graph
may have none, or only inert position-local ones). Let $\iota(v) = \softmax(\Sfac_{v,\cdot})$
and $\Delta(v) = \iota(v)_1 - \iota(v)_2$, with the $\iota(v)$ and the $\Delta(v)$ pairwise
distinct over the finite vocabulary, which is generic in $\Sfac$. With $\lambda = 1$ put
%
\begin{equation}
\label{eq:sharp-energy}
  E(\Q;w) = \psi(\Q_1)\,\varphi(\Q_2)
    - \textstyle\sum_{i=1}^{2}\langle \Sfac_{w_i,\cdot}, \Q_i\rangle,
\end{equation}
%
where $\psi(\Q_1) = h(\Q_1(1) - \Q_1(2))$, $\varphi(\Q_2) = \Q_2(1)$, and $h$ is the univariate
Hermite interpolation polynomial with $h(\Delta(v)) = c_v$ pairwise distinct and
$h'(\Delta(v)) = 0$ at every $v$ in the vocabulary; such an $h$ exists, of degree at most
$2|\Vocab|-1$. Then $E$ is polynomial, hence $C^2$, and its pair term is word-free: SE holds.
C1: whenever $z_1$ is updated its centred field is
$\Proj{z_1}[-\Sfac_{w_1,\cdot} + \varphi(\Q_2)\nabla_{\Q_1}\psi]$, and along the tangent
direction $(1,-1)$ the second term contributes $2h'(\Delta(w_1)) = 0$ at the current state, so
the update returns $\softmax(\Sfac_{w_1,\cdot}) = \iota(w_1)$. The reachable set of $z_1$ is
that fixed point under every schedule, and $\Q_1^{(t)}$ never sees $w_2$. X: under the schedule
$B_1 = \{z_2\}$ the field driving $z_2$ contains
$-\psi(\Q_1^{(0)})\Proj{z_2}(1,0)^{\top} = -c_{w_1}\Proj{z_2}(1,0)^{\top}$ with the $c_v$
pairwise distinct, so $\Q_2^{(1)}$ varies with $w_1$, from past to future, the useful direction.
C2 fails, as it must: off the finite critical set of $h$ we have $h' \neq 0$ and the
cross-position edge fires. The anti-causal coupling exists in $E$ but is never reached.
\end{proof}

C1 holds there only because $h$ is fitted to the rows of $\Sfac$, and perturbing $\Sfac$ with
$h$ fixed breaks it. C2 is what masked attention satisfies for all parameters, activations and
inputs.

\paragraph{Setup.}
$\Nodes$ has a node $z_i$ per label variable $\Zv{i}$ and a node $h_{ic}$ per head variable
$\Hv{i}{c}$, with $\pi(x)$ the position of $x$. Node $x$ carries a distribution $\Q_x$ on a
finite domain of size $d_x$ ($=\nlab$ for label nodes; for a head node the number of candidate
heads, which the theorem leaves arbitrary, not necessarily causal). Write
$\Xsp = \prod_x \Delta_x$ and $\Xint = \prod_x \Delta^\circ_x$,
$\Tang{x} = \{u \in \R^{d_x} : \mathbf{1}^{\top}u = 0\}$, and
$\Proj{x} = I - d_x^{-1}\mathbf{1}\mathbf{1}^{\top}$ for the orthogonal centring projector onto
$\Tang{x}$. A function on $\Xint$ \emph{does not depend on the $y$-block} if it is constant under
changes of the $\Q_y$ coordinates alone. Four assumptions hold throughout.
(A1)~$E \in C^2(U)$ for some open convex $U \supseteq \Xsp$. (A2)~A schedule is a sequence
$(B_1, B_2, \dots)$ of non-empty subsets of $\Nodes$; at step $t$ each $x \in B_t$ is replaced
by $\Q_x^{(t)} = \softmax\big(-\lambda_{x,t}^{-1}\nabla_x E(\Q^{(t-1)})\big)$,
$\lambda_{x,t}>0$, and the rest carry over. This covers synchronous, asynchronous and mixed
schedules and per-node, per-step message weights. (A3)~$E(\Q;w) = E_0(\Q) + \sum_i
u_i(\Q_{z_i}; w_i)$ with $E_0$ word-free and $u_i$ a function of the block $\Q_{z_i}$ alone (in
the encoder, $u_i = -\langle\Sfac_{w_i,\cdot},\Q_i\rangle$), and nothing else depends on $w$.
(A4)~$\Q_{z_i}^{(0)} = \iota_i(w_i) \in \Delta^\circ$, with head initialisations word-independent
and interior. SE is (A1)--(A3) for one common $E$. X is the existence of an $i$, a $t$, a
schedule and words differing only at some coordinate $k \neq i$ for which $\Q_{z_i}^{(t)}$
differs.

\emph{Interiority}: every iterate lies in $\Xint$, since initialisations are interior by (A4) and
softmax outputs are strictly positive, so every gradient below is evaluated at an interior point.
(A1) is needed and is not a technicality. The argument runs through equality of the mixed
partials of $E$, which their existence alone does not give: for
$f(x,y) = xy(x^2-y^2)/(x^2+y^2)$ with $f(0,0)=0$ both exist at the origin and disagree.

\paragraph{Gauge reduction.}

\begin{lemma}
\label{app:lem-gauge}
Fix a node $x$ and $\lambda>0$. For $P, P' \in \Xint$ the updates
$\softmax(-\nabla_x E(P)/\lambda)$ and $\softmax(-\nabla_x E(P')/\lambda)$ agree if and only if
$\Proj{x}\nabla_x E(P) = \Proj{x}\nabla_x E(P')$; so for $y \neq x$ the $x$-update depends on the
$y$-block somewhere on $\Xint$ if and only if the centred field $\Proj{x}\nabla_x E$ does.
\end{lemma}

\begin{proof}
$\softmax(g) = \softmax(g + c\mathbf{1})$ for every $c \in \R$, and $\softmax$ restricted to
$\Tang{x}$ is injective: $\softmax(g) = \softmax(g')$ forces $g - g' = c\mathbf{1}$, which
centring forces to zero. Multiplication by $-1/\lambda$ is a bijection of $\Tang{x}$.
\end{proof}

An update therefore reads only the tangential part of the gradient. This closes the
shift-invariance loophole of the softmax by definition: a coupling invisible to $x$'s softmax is
one that, on the feasible set, does not depend on the $x$-block at all. Write $x \rhd y$ when
$\Proj{x}\nabla_x E$ depends on the $y$-block somewhere on $\Xint$, and let $\Gamma$ be the graph
on $\Nodes$ with an edge $\{x,y\}$ exactly when $x \rhd y$.

\begin{lemma}[No one-way coupling]
\label{app:lem-symmetry}
For $x \neq y$, the conditions $x \rhd y$; $u^{\top}\nabla^2 E(Q)v \neq 0$ for some
$u \in \Tang{x}$, $v \in \Tang{y}$ and $Q \in \Xint$; and $y \rhd x$ are equivalent. In
particular $\rhd$ is symmetric and $\Gamma$ is a well-defined undirected graph.
\end{lemma}

\begin{proof}
If $u^{\top}\nabla^2 E(Q)v \neq 0$ then $Q+sv \in \Xint$ for $|s|$ small, one block moving inside
its open simplex, and $g(s) = u^{\top}\Proj{x}\nabla_x E(Q+sv) = u^{\top}\nabla_x E(Q+sv)$ has
$g'(0) \neq 0$. It is non-constant along a segment whose endpoints differ only in the $y$-block,
so $x \rhd y$. Conversely, if that tangential Hessian block vanishes throughout $\Xint$, take
$P,P' \in \Xint$ differing only in the $y$-block. The segment between them lies in $\Xint$ by
convexity of $\Delta^\circ_y$ and the product structure, and
$\phi(s) = \Proj{x}\nabla_x E(P + s(P'{-}P))$ is $\Tang{x}$-valued with $\phi'(s)$ orthogonal to
$\Tang{x}$. Since $\Tang{x}$ is closed, $\phi' \in \Tang{x} \cap \Tang{x}^{\perp} = \{0\}$, so
$\phi$ is constant and $\neg(x \rhd y)$. The middle condition is symmetric under
$(x,u) \leftrightarrow (y,v)$ because $\nabla^2 E(Q)$ is a symmetric matrix at every $Q$
(Schwarz's theorem; this is where (A1) is used).
\end{proof}

\paragraph{Causality removes both directions.}
Write $C(x)$ for the $\Gamma$-connected component of $x$. Induction over an arbitrary schedule
gives \emph{confinement}: $\Q_x^{(t)}$ is a function of the initial blocks
$(\Q_y^{(0)})_{y \in C(x)}$ and the words $(w_k)_{z_k \in C(x)}$ alone. Indeed,
\Cref{app:lem-gauge} makes the update of $x \in B_t$ a function of $\Proj{x}\nabla_x E$ and, for
a label node, of $w_{\pi(x)}$ and of no other word by (A3). \Cref{app:lem-symmetry}, applied one
block at a time inside $\Xint$, leaves that field a function of the blocks of $x$ and its
$\Gamma$-neighbours, all in $C(x)$. Head nodes need no separate treatment, since a route
``through an $H$'' is a sequence of $\Gamma$-edges incident to that node. Confinement is a
necessity statement (no path implies no influence), not a transmission statement: couplings can
cancel or sit dead on the reachable set, which is what \Cref{app:prop-sharp} exploits.

\begin{proposition}
\label{app:prop-c2}
Assume SE. Then C2 holds if and only if $\Gamma$ has no edge between nodes at different
positions.
\end{proposition}

\begin{proof}
Let $\pi(x) < \pi(y)$ and apply C2 to the one-node schedule $B_1 = \{x\}$ at $t=1$: for every
$\Q^{(0)} \in \Xint$ the update of $x$ ignores the $y$-block. Because C2 quantifies over all
interior initialisations, this is an identical statement on $\Xint$, so $\neg(x \rhd y)$ by
\Cref{app:lem-gauge} and then $\neg(y \rhd x)$ by \Cref{app:lem-symmetry}. This is the crux of
the theorem: Schwarz symmetry converts ``the later node's block does not feed the earlier node's
update'' into ``the earlier node's block does not feed the later node's update''. The
cross-position pair was arbitrary. Conversely, position-pure components make every $\Q_x^{(t)}$
a function of position-$\pi(x)$ initial blocks and of $w_{\pi(x)}$ only, by confinement, which
is C2.
\end{proof}

\begin{proof}[Proof of \Cref{thm:trilemma}]
Under SE and C2, \Cref{app:prop-c2} leaves $\Gamma$ without a cross-position edge, so every
component lies inside one position's node set. Word entry is confined to label nodes by
(A3)--(A4), and confinement then makes $\Q_{z_i}^{(t)}$ a function of $w_i$ and the
position-$i$ initial blocks alone, for every schedule, every $t$ and every $i$. No label marginal
ever depends on another word, which is $\neg$X. Contrapositively, a network-causal contextual
scheme is not generated, in the sense of (A2), by any single $C^2$ word-local energy. C2
$\Rightarrow$ C1 follows as well, word initialisations being particular interior ones; it holds
vacuously, since nothing depends on any other word. Under a single energy the only way to be
causal is to be trivial. The other half of the statement, that any two of the three properties
hold together, is given by the three witnesses below.
\end{proof}

For the masked scheme one can say more. One tangential Hessian entry is computed from the update
at $i$ and again from the update at $j<i$, whose masked message contains no $\Q_{z_i}$ at all.
The two must agree, which forces $\Proj{}\Tfac{c}\Proj{} = 0$ on every channel, and the centred
label message of such a $\Tfac{c}$ is a state-independent constant. A non-degenerate
$\Tfac{c}$ admits no energy, and a degenerate one carries no context. This is the formal content
of \Cref{rmk:never-formed}, and the reason the third witness below must give up SE.

\paragraph{The three witnesses.}
All three pairs are realisable, so the impossibility in \Cref{thm:trilemma} needs all three
properties. \emph{SE $\wedge$ C2, $\neg$X:} take
$E = \sum_i u_i(\Q_i;w_i) + \sum_{i,c}\psi_{ic}(\Q_{z_i},\Q_{h_{ic}})$, whose couplings are all
within one position, so C2 holds by \Cref{app:prop-c2}. \emph{SE $\wedge$ X, $\neg$C2:} the
encoder model of \citet{wu2023probabilistic}, whose trilinear term
$\Q_{h_{jc}}(i)\,\Q_j(b)\,\Q_i(a)\,\Tfac{c}_{b,a}$ makes a cross-position edge that
\Cref{app:lem-symmetry} fires in both directions. \emph{C2 $\wedge$ X, $\neg$SE:} mask the head
domains so that $\Hv{i}{c}$ may point only to positions $j<i$, and keep only the message term
weighted by $\Q_{h_{ic}}(j)$. Every activation at position $i$ is then a function of blocks at
positions $\leq i$ for arbitrary interior states, X holds for generic $\Tfac{c}$, and $\neg$SE is
forced by \Cref{thm:trilemma}.

\paragraph{Back to the model.}
The directed chain of \Cref{sec:chain} gives up the single global energy and keeps causality and
contextuality. It is not the third witness: that scheme deletes a term from a symmetric field,
whereas the chain never forms one (\Cref{rmk:never-formed}). The energy is not simply dropped
either. Each step conditional is generated by its own energy $E_i$, every update is exactly
$\softmax(-\nabla E_i/\lambda)$ with nothing deleted, and the global object is the valid directed
joint $\prod_i p(z_i, h_i \mid z_{\prefix{i}}, w_{1:i})$. Inference is variational filtering, so
the belief at $i$ approximates $p(\Zv{i} \mid w_{1:i})$ and not the smoothing marginal
$p(\Zv{i} \mid w_{1:\seqlen})$. For a decoder, filtering is the right object.
